%-*-tex-*-
\ifundefined{writestatus} \input status \relax \fi %
\chcode{break}

\def\cqu{}

\chapterhead{break}{BREAKS\cr and\cr NOBREAKS}
\tex\ controls the breaking of lines and pages by the judicious insertion of
penalties, both positive and negative. A negative penalty will encourage
a break
and a positive penalty will inhibit it. The commands in this chapter either
inhibit or cause breaks. In addition, most of them have some form of vertical
skip. All but |\linebreak| are intended to be used in vertical mode. In all
cases, the use of a |\...break| in horizontal mode will first cause a change
to (internal) vertical mode where the penalty will be applied. A penalty of
$\pm$10000 is the same as $\pm$ infinity. 

\shead{breakcoms}{Command Forms}
Since the commands generally do not take parameters, the command list is
being dispensed with.
\bshortcomlist
\pla\@|\allowbreak|&no penalty for a break \dots\ |\penalty 0|\cr
\pla\@|\break|&causes a break \dots\ |\penalty 10000|\cr
\pla\@|\bigbreak|&encourages break and vskips the maximum of the
|\bigskipamount| and the |\lastskip| \dots\ |\penalty -200| \dots\ skip is 
\quad\hbox to .25in{\vbox{\hrule width .25in\bigbreak\hrule}} \cr
\pla\@|\bye|&terminates a document and vertical |fill|s the page\cr
\pri|\discretionary|&a discretionary break or hyphen. The form is
\@|\discretionary{<pre  break  text>}{<post  break  text>}{<text>}|. If
there is a break, the |<pre break text>| goes on the first line and the
|<post break text>| goes on the second line. If there is nobreak, the
|<text>| is used.\cr
\ext\@|\done|&terminates a document and vertical |fill|s the page and does
some tests on groups used in \intex. \cr
\ext\@|\end|&terminates a document without doing any fill on the last page\cr
\pla\@|\filbreak|&encourages a break with |\penalty -200| and vertical |fils|
the page. It is useful for address lists where it is desirable not to split
an individual address at a page boundary\cr
\pla\@|\eject|&causes a column or page break but does not vertical |fill| the
page. Its use usually causes underfull boxes.\cr
\ext\@|\ejectcolumn|&causes a column break (equivalent to a page break in single
column) and vertical |fills| the page\cr
\ext\@|\ejectpage|&causes a page break and vertical |fills| the page \dots\ in
multicolumn this will cause a new page, not just a new column\cr
\pla\@|\goodbreak|&encourages a break \dots\ |\penalty -200|\cr
\ext\@|\linebreak|&forces a line break in horizontal mode {\bf and} |fils| the
    line forcing the words to the left side of the line\cr
\pla\@|\medbreak|&encourages break and vskips the maximum of the
         |\medskipamount| and the |\lastskip| \dots\ |\penalty -100| \dots\ skip is 
         \hbox to .25in{\vbox{\hrule width.25in\medbreak\hrule}} \cr
\pla\@|\nobreak|&does not allow a break \dots\ |\penalty 10000|\cr
\ext|\penaltybreak|&form \@|\penaltybreak{<penalty>}{<skip size>}| \dots\ 
          this is a generalization of |\medbreak| allowing arbitrary penalties and skip
          sizes\cr
\pla\@|\smallbreak|&encourages break and vskips the maximum of the
|\smallskipamount| and the |\lastskip| \dots\ |\penalty -50| \dots\ skip is 
\hbox to .25in{\vbox{\hrule width .25in\smallbreak\hrule}} \cr
\eshortcomlist
\ejectpage
\done